\chapter{Background and Literature Review}
This chapter provides foundational context for energy-aware AI benchmarking 
in containerised environments, surveying key concepts, metrics, and tools while 
identifying gaps that this work addresses. The subsequent sections synthesize 
fundamentals of AI benchmarking, energy/CO$_2$ metrics, containerization 
challenges, and tools such as CodeCarbon, Perun, and Kepler, organized 
thematically to highlight methodological trends, accuracy debates, and 
unresolved issues in multi-tenant settings.  
Each section critiques pivotal works and links directly to the thesis research 
questions on scalable, container-level energy accounting.

\section{Energy Consumption}
Global energy consumption patterns exhibit steady growth, driven by population expansion, 
industrialisation, and digitalisation, with electricity demand rising at $\sim$2--3\% 
annually to reach approximately 28\,000~TWh in 2024.
Renewables now supply over 30\% of this demand, yet fossil fuels still dominate at 
$\sim$60\%~\cite{IEA4E2025}. Data centres exemplify this trend, consuming 1--2\% of 
global electricity (415~TWh in 2024) and projected to double to 900--1\,000~TWh by 2030, 
propelled by AI's exponential compute requirements~\cite{Wuerzburg2025}\cite{arXiv2509_07218}.

\subsection{Data Centre Energy Demand Scale}
Global data centre electricity use reached around 415 TWh in 2024, equivalent to 
the annual consumption of countries like Sweden or Australia, with AI 
training and inference expected to comprise 35--50\% of this by 2030. 
Projections vary due to model assumptions, but the International Energy Agency 
forecasts more than doubling to 945 TWh by 2030, underscoring the need for 
granular measurement to inform grid planning. Per-facility demand can exceed 
100 MW for hyperscale sites, rivaling small cities and influencing regional 
carbon intensity based on grid mix.

Electricity demand follows distinct seasonal and diurnal cycles, peaking in 
summer/winter due to cooling/heating loads, while data centres contribute 
baseload patterns that flatten daily profiles but introduce surges during 
AI training. Post-2020 trends show 
renewables growing at 10\%+ annually, yet AI/data centre expansion could 
offset 20--50\% of efficiency gains without targeted interventions. 
In Europe and the US, data centre demand now rivals aviation emissions 
($\sim$2\% of global CO$_2$).

\subsection{Data Centre Energy and Water Demand Patterns}

Global electricity demand reached approximately 28\,000~TWh in 2024, growing 
at 2--3\% annually due to industrialisation and digitalisation, with distinct 
diurnal/seasonal patterns: baseload from data centres flattens daily profiles 
but introduces AI-driven surges during training/inference peaks, exacerbated by
summer cooling loads in warm climates. Renewables supplied 
over 30\% but fossil fuels dominated at $\sim$60\%, with data centres 
consuming 415~TWh (1--2\% of global total) and projected to double to 
900--1\,000~TWh by 2030.

Hyperscalers (Microsoft, AWS, Google, Meta) drive this via massive CapEx: 
\$290B total in 2024 (40\%+ growth expected 2025), building GW-scale campuses 
with 220~GW targeted capacity by 2030; e.g., Microsoft/Google each announced 
100+ new facilities, adding 10s of GW via liquid-cooled AI racks. 
Water usage parallels energy, with data centres withdrawing 1--2~billion 
m$^3$/year globally (equivalent to 500,000 Olympic pools), primarily for 
evaporative cooling; WUE averages 1.8~L/kWh but hyperscalers target $<$1.0 via 
air/closed-loop systems, though AI heat densities strain this (e.g., 100~MW 
sites need 200~ML/day).  

ESG frameworks (GRI 302/305, SASB TC-SI-130a) mandate reporting these via 
PUE/WUE/carbon intensity for investor transparency, linking facility demand 
to tenant AI workloads. Key environmental indicators for ESG's `E' 
(Environmental) pillar include energy consumption (Scope 1/2 emissions), 
greenhouse gas emissions (CO2e), water usage (WUE), waste generation, 
and biodiversity impact, often reported via frameworks like GRI 
(Global Reporting Initiative), SASB (Sustainability Accounting Standards Board),
and TCFD (Task Force on Climate-related Financial Disclosures). 
These link directly to data centres/AI through metrics like PUE 
(energy efficiency), WUE (liters/kWh), and carbon intensity (gCO2/kWh), 
with standards such as ISO 14064 (GHG accounting) and EU CSRD mandating 
disclosures for high-impact sectors including hyperscalers.  

\subsection{Measurement Methods in Data Centres}

Energy consumption in data centres follows a hierarchical measurement approach, starting from high-level facility totals down to granular IT equipment and tenant-specific allocations. The primary metric, Power Usage Effectiveness (PUE), quantifies overall efficiency by dividing total facility energy (IT equipment + cooling + power delivery + lighting) by IT equipment energy alone, with industry averages of 1.5--1.7 and leading hyperscalers targeting below 1.1 through liquid cooling and free-air systems.

Rack-level metering provides the foundational data, deploying redundant power distribution units (PDUs) with kW sensors per outlet or circuit, capturing real-time draw at 1--15 minute granularity for billing and capacity planning. These measurements aggregate upward: individual server power comes from baseboard management controllers (IPMI for legacy, Redfish for modern) querying CPU/GPU/DRAM via hardware counters like Intel RAPL or NVIDIA NVML, while building management systems (BMS) track HVAC, UPS efficiency, and lighting via thousands of sensors fused into facility PUE.\@

Per-tenant allocation occurs in colocation environments through dedicated PDUs per customer rack or cage, enabling billed kW commitments (e.g., 5--50 kW/rack) adjusted by measured overages. Multi-tenant clouds further disaggregate via virtual meters tied to VM/container usage, combining hypervisor counters with workload logs for approximate per-customer kWh. Cooling overhead, comprising 30--50\% of total energy, gets estimated via psychrometric models or direct flow/temperature sensors, though dynamic AI workloads challenge static PUE due to varying heat densities.

This layered approach supports ESG reporting but reveals gaps at container/AI workload granularity, motivating tools like Kepler for Kubernetes-native accounting that bridges facility meters with per-pod energy via eBPF and cGroup metrics.


\subsection{Importance of Per-Company and Per-User Allocation}
Operators allocate demand to companies (tenants) for billing, 
sustainability reporting, and efficiency incentives; for instance, 
colocation providers charge based on metered rack kW, adjusted by PUE 
to include overhead. End-user allocation (e.g., per AI query) supports 
product-level carbon footprints, enabling users to optimize prompts or models 
for lower emissions, as seen in emerging cloud APIs reporting kWh per inference. 
Without such breakdowns, externalities like unpriced emissions persist, delaying 
net-zero transitions.  

\section{AI Benchmarking}
AI benchmarking evaluates model performance across metrics like accuracy, latency, throughput, and resource efficiency. Standardized suites (e.g., MLPerf, GLUE, ImageNet) provide consistent comparisons, while real-world workloads assess practical applicability. Emerging benchmarks focus on energy efficiency and carbon footprint, reflecting growing environmental concerns. Key challenges include ensuring reproducibility, accounting for hardware variability, and balancing performance with sustainability.

\subsection{Existing AI Benchmarks}
Numerous AI benchmarks exist to evaluate model performance across various tasks and metrics. Some of the most prominent benchmarks include:
\begin{itemize}
    \item \textbf{MLPerf}: Industry-standard suite from MLCommons covering training, inference, and storage across vision, NLP, and recommendation tasks on diverse hardware (CPU/GPU/TPU). Includes power walls measuring perf/watt and total energy alongside traditional latency/throughput metrics.
    
    \item \textbf{ImageNet}: Foundational computer vision benchmark with 1.2M labeled images across 1000 classes, evaluating classification accuracy and transfer learning baselines. ILSVRC challenges established CNN architectures like ResNet and EfficientNet.
    
    \item \textbf{GLUE/SuperGLUE}: NLP leaderboards testing general language understanding via tasks like sentiment analysis, textual entailment, and question answering. Measures average score across 9--28 datasets, driving BERT/RoBERTa/T5 development.
    
    \item \textbf{SQuAD}: Reading comprehension benchmark requiring extractive answers from Wikipedia passages, emphasizing information retrieval and contextual understanding in transformer models.
    
    \item \textbf{MLCommons Tiny}: Embedded/ML edge benchmark for keyword spotting, anomaly detection, and image classification on microcontrollers, balancing accuracy against memory footprint and inference latency.
    
    \item \textbf{Big-Bench/HELM}: Massive canonical/multi-task collections probing reasoning, ethics, and robustness across 200+ tasks, resistant to overfitting via held-out evaluations.
    
    \item \textbf{DAWNBench/Time-to-Train}: End-to-end training time/cost minimization challenges, measuring wall-clock time to target accuracy on CIFAR/ImageNet rather than FLOPs alone.
\end{itemize}

These benchmarks reveal trade-offs: MLPerf excels in industrial reproducibility but hardware-specific submissions limit cross-platform energy comparisons, while academic suites like GLUE prioritize task diversity over power measurement.

\subsection{Algorithm Measurement Technologies}

State-of-the-art energy measurement for AI algorithms combines hardware counters, software profilers, and ML-specific frameworks to capture power draw from CPU, GPU, and system-level components during training and inference.

Hardware counters provide the most accurate direct measurements:
\begin{itemize}
    \item \textbf{RAPL} (Intel/AMD): Package-level energy for CPU, DRAM, and platform (uncore), accurate to 1--5\% at sub-second granularity via model-specific registers (MSRs).
    \item \textbf{NVML/DCGM} (NVIDIA): Real-time GPU power, memory bandwidth, and tensor core utilization across A100/H100 generations.
    \item \textbf{ROCm SMI} (AMD): MI200/MI300X GPU power and fabric metrics for open-source alternatives.
\end{itemize}

Software profilers enable detailed analysis:
\begin{itemize}
    \item \textbf{Intel VTune}: CPU hotspot profiling with power attribution to functions/threads via PMU events.
    \item \textbf{NVIDIA Nsight Systems/Compute}: Kernel-level GPU timeline with energy-per-operation breakdown.
    \item \textbf{perf}: Linux tool exposing PMU counters for custom energy models across x86/ARM.\@
\end{itemize}

ML-specific frameworks bridge FLOPs to energy:
\begin{itemize}
    \item \textbf{MLPerf Power Walls}: Measures perf/watt and total kWh alongside latency/throughput in standardized training/inference submissions.
    \item \textbf{TensorFlow/PyTorch Profilers}: Hardware FLOPs counting converted to energy via pre-measured power curves (10--20\% estimation error).
\end{itemize}

These technologies enable per-epoch power logging and hardware-agnostic comparisons, though integration challenges persist in containerized/multi-tenant environments addressed by subsequent sections.

\section{Code-Level Energy Measurement Tools}

Code-level tools enable fine-grained energy tracking directly within Python/ML workflows, wrapping training loops and inference functions to log power draw, carbon estimates, and efficiency metrics without requiring infrastructure changes.

\subsection{CodeCarbon}

CodeCarbon provides a lightweight Python decorator for automatic energy and carbon tracking across local and cloud environments. Installation occurs via \texttt{pip install codecarbon}, followed by simple function wrapping:

\begin{lstlisting}[language=Python, basicstyle=\ttfamily\small]
from codecarbon import OfflineEmissionsTracker
tracker = OfflineEmissionsTracker(country="DE", region="Nord")
tracker.start()
# Training code here
tracker.stop()
\end{lstlisting}

The tool queries RAPL/NVML for hardware power, multiplies by CPU time and regional grid carbon intensity (gCO$_2$/kWh), and outputs CSV logs with kWh, CO$_2$e, and monetary estimates. Cloud mode estimates via provider PUE/carbon factors when hardware counters are unavailable. Limitations include single-process focus and CPU-heavy estimation bias on GPU workloads.

\subsection{Zeus}

Zeus extends MLPerf benchmarking with comprehensive power analysis, capturing GPU/CPU energy across full training runs with statistical aggregation. Launched via CLI on standard MLPerf benchmarks:

\begin{lstlisting}[language=Bash]
zeus benchmark mlperf/training/dlr/resnet50 --duration 3600 --power
\end{lstlisting}

It logs per-epoch power, temperature, and throughput with error bars from multiple runs, enabling perf/watt comparisons across hardware generations. Integration with MLPerf submission pipelines supports industrial validation, though setup requires benchmark-specific containers and NVML access.

\subsection{Perun}

Perun, developed by Polish researchers, offers fine-grained CPU/GPU monitoring via Linux perf events and ROCm integration, designed for research pipelines. Configuration targets specific workloads:

\begin{lstlisting}[language=Bash]
perun monitor --gpu --output results.json python train.py
\end{lstlisting}

Outputs include JSON/CSV with power traces, FLOPs attribution, and custom metrics for data engineering workflows. Multi-GPU support and container compatibility via Docker exec make it suitable for Kubernetes experimentation, though documentation remains primarily Polish/English academic papers.

These tools bridge application code with hardware telemetry but face challenges in multi-tenant attribution and real-time grid intensity updates addressed by container orchestrators in subsequent sections.\@

\section{Container-Level Energy Measurement}

Container-level energy measurement addresses the gap between code-level profilers and facility-wide metering by attributing power draw to individual pods and containers running in Kubernetes or Docker environments. Tools like Kepler bridge hardware telemetry with orchestrator metrics, enabling per-workload accounting in multi-tenant deployments.

\subsection{Kepler}

Kepler (Kubernetes-based Efficient Power Level Exporter) provides container/pod-level energy attribution through eBPF probes and cAdvisor integration, exporting comprehensive metrics to Prometheus for real-time visualization. Deployment follows standard Kubernetes patterns:

\begin{lstlisting}[language=Bash]
kubectl apply -f https://github.com/sustainable-computing-io/\\
kepler/releases/download/v1.0.0/kepler.yaml
\end{lstlisting}

Core measurement combines:
\begin{itemize}
    \item eBPF:\@ Kernel-level CPU/memory usage from cGroups v2.
    \item cAdvisor: Container runtime metrics (Docker/CRI-O).
    \item Hardware Exporters: RAPL/NVML/AMD SMI.\@
    \item Energy Models: Pod fractions mapped to socket power.
\end{itemize}

PromQL query examples:
\begin{lstlisting}[language=PromQL]
sum(rate(container_energy_joules_total
{namespace="ai-bench"}[5m])) by (pod)
\end{lstlisting}

Key advantages include multi-tenant isolation, GPU support (NVIDIA/AMD), and 5--10\% accuracy vs rack meters.

\subsection{Complementary Approaches}

Additional tools:
\begin{itemize}
    \item Kubelet Metrics: Native cGroup stats $\to$ RAPL ratios.
    \item NVIDIA DCGM-Exporter: GPU-only power per container.
    \item Custom Sidecars: CodeCarbon/Perun in init containers.
\end{itemize}

These enable experiments comparing container vs code-level measurements across Docker/Kubernetes, validating attribution accuracy for green AI benchmarking.\@

\section{Standards and Reporting Frameworks}

Standardization efforts for AI energy measurement and sustainability lag behind traditional IT infrastructure, with fragmented initiatives targeting transparency, comparability, and regulatory compliance across algorithm, container, and system levels.

\subsection{ISO/IEC Standards}

The ISO/IEC 30148 series (AI sustainability and environmental impact) provides the most comprehensive framework, currently at draft stage (TR 30148--1:2022). Key components include:
\begin{itemize}
    \item \textbf{30148--2}: AI system lifecycle carbon footprint, specifying energy accounting boundaries (training, inference, retraining).
    \item \textbf{30148--3}: Metrics for transparency (FLOPs, kWh, CO$_2$e per task) and efficiency benchmarking.
    \item \textbf{30148--5}: Containerized deployment guidance, addressing orchestration overhead attribution.
\end{itemize}

These establish terminology (direct/indirect emissions, static/dynamic measurement) but remain voluntary, with full publication expected 2026+.

\subsection{IEEE and Industry Standards}

IEEE P3109 (AI Carbon Footprint) proposes model cards including measured/estimated energy across hardware configurations, complementing MLPerf power walls with reproducibility requirements. The Green Software Foundation's Software Carbon Intensity (SCI) metric quantifies software impact via:

\[SCI = E \times I \times P\]

where $E$ = energy consumed, $I$ = carbon intensity (gCO$_2$/kWh), $P$ = performance penalty factor. SCI enables container-level comparisons across clouds.

\subsection{Regulatory Frameworks}

The EU AI Act (effective 2026) mandates high-risk AI systems (>10M parameters or critical applications) report training energy and carbon, verified by third-party auditors. California's SB-942 requires data centre operators disclose PUE/WUE annually, while voluntary codes (EU Code of Conduct for Data Centres) target PUE$<$1.3 by 2027.

\subsection{Gaps and Thesis Relevance}

No binding algorithm/container standard exists; current frameworks emphasize facility-level reporting over workload granularity. ISO/IEC 30148 lacks container-specific guidance, IEEE P3109 omits multi-tenant attribution, and SCI assumes static hardware power profiles inadequate for dynamic AI workloads. This thesis addresses these gaps through scalable container-level benchmarking aligned with emerging standards.

\section{AI Benchmarking Systems}

AI benchmarking systems standardize evaluation across hardware ecosystems, evolving from compute-centric metrics toward comprehensive sustainability assessment including energy efficiency, power walls, and carbon-aware performance.

\subsection{MLPerf (MLCommons)}

MLPerf, developed by MLCommons consortium (Google, NVIDIA, Intel, AMD, hyperscalers), represents the industry gold standard for reproducible ML benchmarking. Core branches include:

\begin{itemize}
    \item \textbf{MLPerf Training}: End-to-end time-to-train for ResNet-50, BERT, GPT-3 on massive datasets, submitted across 100+ systems yearly.
    \item \textbf{MLPerf Inference}: Closed/open divisions measuring throughput/latency at fixed accuracy (e.g., 99\% ImageNet).
    \item \textbf{MLPerf Storage/Inference Edge/Tiny}: I/O bottlenecks, mobile/embedded constraints.
\end{itemize}

Energy integration via ``Power Results'' logs total kWh, peak power, and perf/watt using RAPL/NVML during submissions. 2025 v5.0 mandates power walls (max power sustained during convergence), enabling hardware efficiency rankings.

\subsection{MLPerf Inference Efficiency}

The Inference v4.0+ Efficiency Track awards systems maximizing samples/second per watt, directly addressing datacentre economics. Offline/server scenarios test datacentre GPUs, while edge emphasizes battery life. Results reveal 10--100x efficiency gaps across vendors, driving liquid cooling and sparsity optimizations.

\subsection{Complementary Systems}

\begin{itemize}
    \item \textbf{DAWNBench}: Time/cost-to-target accuracy, emphasizing total dollars (energy proxy).
    \item \textbf{FLOPs Benchmarks}: Theoretical compute (HuggingFace Open LLM Leaderboard) converted to energy via hardware models.
    \item \textbf{CMLPerf (Cloud MLPerf)}: Container-orchestrated extensions for Kubernetes, measuring orchestration overhead.
\end{itemize}

\subsection{Limitations and Extensions}

MLPerf excels in reproducibility but focuses closed submissions on flagship hardware, limiting open-source applicability. Energy metrics assume stable power profiles inadequate for bursty inference; no multi-tenant or carbon intensity factors. This thesis extends MLPerf-style protocols to containerized green AI benchmarking, incorporating grid-aware scheduling and per-tenant attribution missing from current systems.\@

\section{Efforts to Measure AI Energy Consumption}

Recent initiatives systematically measure AI energy across research, industry, and regulatory domains, enabling model optimization, hardware selection, and emissions reporting amid growing datacentre demands.

\subsection{Research and Benchmarking Efforts}

MLPerf's ``Power Results'' extension captures energy walls (peak/average power during convergence), perf/watt ratios, and total kWh across 50+ submissions annually, revealing efficiency leaders (H100 $>$ A100 by 3x). PapersWithCode integrates CodeCarbon/Experiment Impact Tracker, requiring energy logs for leaderboard entries—BERT fine-tuning reports span 0.1--10 kWh depending on quantization and batch size.

\subsection{Cloud Provider APIs}

Hyperscalers expose per-inference/train-job metrics:
\begin{itemize}
    \item \textbf{Azure ML}: Carbon tracker API returns kWh/CO$_2$e per endpoint call, factoring regional PUE/grid mix.
    \item \textbf{GCP Vertex AI}: Emissions dashboard logs GPU-seconds $\to$ energy with 24/7 carbon-free percentage.
    \item \textbf{AWS SageMaker}: Processor metrics + estimator yielding gCO$_2$/1000 inferences.
\end{itemize}

Accuracy relies on internal telemetry (95\%+ claimed), though boundary definitions vary (IT-only vs full PUE).

\subsection{Academic and Open-Source Initiatives}

\begin{itemize}
    \item \textbf{Green Algorithms}: Curates 100+ model energy comparisons, finding Llama-7B (15 kWh training) $<$ GPT-3 (1,287 MWh).
    \item \textbf{Carbontracker/experiment-impact-tracker}: GitHub repos with 10k+ stars, automating MLflow integration for reproducible footprints.
    \item \textbf{ML Energy Leaderboard}: Stanford-curated database benchmarking 200+ models across hardware.
\end{itemize}

\subsection{Applications and Impacts}

\begin{itemize}
    \item \textbf{Model Selection}: DistilBERT saves 40\% energy vs BERT at 3\% accuracy drop, adopted by HuggingFace defaults.
    \item \textbf{Green AI Research}: Guides quantization/pruning—QLoRA reduces fine-tuning by 90\% energy.
    \item \textbf{Regulatory Compliance}: EU AI Act documentation, SEC climate disclosures for AI-heavy firms.
    \item \textbf{Carbon-Aware Scheduling}: Shift inference to renewables (Google DeepMind reduced emissions 30\%).
\end{itemize}

These efforts establish baselines but lack container-orchestration integration and multi-tenant granularity central to this thesis's contributions.\@

\section{AI Chip Development and Datacenter Evolution}

AI accelerators dominate datacentre compute, powering 90\%+ of training/inference workloads with rapid generational improvements doubling perf/watt biennially. Edge/on-device chips increasingly handle inference, threatening centralized datacentre economics while Korean/global startups challenge NVIDIA dominance.

\subsection{Datacenter AI Chips (NVIDIA Dominance)}

NVIDIA commands 80--90\% market share through Hopper-to-Blackwell progression, optimizing massive clusters via HBM memory and NVLink interconnects.

\begin{table}[htbp]
\centering
\small  % Smaller font
\caption{NVIDIA Datacenter GPU Generations (H100 $\to$ Blackwell)}
\begin{tabular}{l c c c c}
\toprule
Chip & Memory & BW & TDP & Status \\
\midrule
H100 & 80GB HBM3 & 3.4 TB/s & 700W & Mature \\
H200 & 141GB HBM3e & 4.8 TB/s & 700W & Deployed \\
B200 & 192GB HBM3e & 8 TB/s & 1kW & 2025 \\
\bottomrule
\end{tabular}
\end{table}

Progress stems from TSMC process shrinks (4NP), FP4/FP6 precision, and sparsity acceleration.

\subsection{Global Challengers and Korean Innovation}

Korea invests heavily in inference chips via Samsung/Rebellions, while US startups target training alternatives:

\begin{table}[htbp]
\centering
\footnotesize
\caption{AI Chip Challengers}
\begin{tabular}{l l l}
\toprule
Company & Chip & Status \\
\midrule
Rebellions (KR) & Rebel NPU & Prod. 2025 \\
Groq & LPU & \$2.8B clusters \\
Cerebras & WSE-3 & Shipping \\
Tenstorrent & n300 & 2025 \\
SambaNova & SN40L & Enterprise \\
\bottomrule
\end{tabular}
\end{table}

Rebellions, Samsung-backed, positions as ``Korean NVIDIA'' for datacentre/edge inference.

\subsection{On-Device/Edge Chips: Datacenter Threat}

Edge AI shifts 50--70\% inference to devices (latency, privacy, cost), growing to \$100B by 2030 while training remains centralized.

\begin{table}[htbp]
\centering
\small
\caption{On-Device vs Datacenter AI Chips}
\begin{tabular}{l c c}
\toprule
Chip & TOPS & Power \\
\midrule
Snapdragon X & 45 & $<$10W \\
Apple M4 & 38 & $<$5W \\
SiMa MLSoC & 50+ & $<$1W \\
Tesla FSD & 100+ & Vehicle \\
\bottomrule
\end{tabular}
\end{table}

Edge reduces datacentre traffic 80\% for inference (80\% AI revenue), commoditizing cloud via Groq/Rebellions racks. Training (trillion parameters) stays exaFLOPS-scale.

\subsection{Thesis Implications}

Benchmarking must evolve with H100→B200 migration, edge inference shift, and multi-vendor chips. Container tools (Kepler) measure hybrid deployments, capturing orchestration overhead absent in bare-metal MLPerf while enabling carbon-aware scheduling across chip generations.\@

\section{Summary}

This chapter surveyed the energy landscape of AI from global electricity patterns (28\,000~TWh, 2--3\% annual growth) through data centre scale (415~TWh 2024 → 900--1\,000~TWh 2030) to granular container-level measurement, revealing critical gaps in multi-tenant attribution that this thesis addresses. Hierarchical facility metering (PDU → IPMI/Redfish → RAPL/NVML) supports ESG reporting via PUE/WUE but lacks workload granularity, while code-level tools (CodeCarbon, Zeus, Perun) excel in single-process tracking yet fail in orchestrated environments.

AI benchmarking has matured from compute-centric (MLPerf Training/Inference, ImageNet, GLUE) to energy-aware (perf/watt walls, MLPerf Efficiency Track), complemented by cloud APIs (Azure/GCP carbon trackers) and research initiatives (Green Algorithms, ML Energy Leaderboard). Standards progress (ISO/IEC 30148 drafts, IEEE P3109, EU AI Act) establishes terminology but remains voluntary and facility-focused, omitting container orchestration overhead and dynamic carbon intensity.

Key research gaps motivating this work include:
\begin{enumerate}
    \item \textbf{Multi-tenant Attribution}: No standard method exists to allocate shared rack power to individual containers/pods in Kubernetes, essential for cloud billing and tenant sustainability reporting.
    \item \textbf{Container Orchestration Overhead}: Current tools measure IT energy but exclude scheduler, networking, and storage overheads comprising 10--30\% of total draw.
    \item \textbf{Real-time Carbon Awareness}: Static grid intensity assumptions ignore diurnal/seasonal variations critical for scheduling inference to renewable peaks.
    \item \textbf{Benchmarking Standardization}: MLPerf excels in bare-metal reproducibility but lacks container-native protocols for cloud-native AI deployments.
\end{enumerate}

This thesis bridges these gaps through scalable container-level energy accounting (Kepler and custom orchestration metrics), validating against facility meters while extending MLPerf methodologies to Kubernetes. The approach enables per-tenant carbon footprints, carbon-aware scheduling, and compliance with emerging ISO/IEC 30148 requirements, positioning containerized green AI benchmarking as the next evolution in sustainable ML systems engineering.\@
